\chapter{Marco teórico}
\label{cap:marco}

Este capítulo desarrolla antecedentes, bases conceptuales y teóricas, mapa semántico, estado del
arte por ejes y brecha de investigación. Las definiciones apoyan el planteamiento del Capítulo~1 sin
sustituir la situación subyacente del §1.4.

\section{Antecedentes de la investigación}
\label{sec:antecedentes}
Los antecedentes de esta investigación provienen de una revisión sistematizada de catorce trabajos
(Fase~0, Grupo~9.2), agrupados en seis ejes (E1--E6) con fichas en el Anexo~B. La síntesis
extendida aparece en las secciones~\ref{sec:estado-arte} y~\ref{sec:brecha}; aquí se condensan las
cuatro líneas que convergen hacia el diseño del modelo de alineación.

\begin{enumerate}
  \item \textbf{Madurez de OpenAPI.} Casas et al.\ (2021), en un mapeo sistemático de 47 artículos
        (2011--2020), confirman que OpenAPI/Swagger es el estándar de facto para describir APIs REST
        y que su uso trasciende la documentación ---generación, pruebas, gobernanza---. Ninguno de
        esos trabajos aborda la coherencia con modelos de negocio bancarios de referencia como BIAN.
  \item \textbf{Correspondencia estructural y semántica.} Shvaiko y Euzenat (2005) consolidan el
        marco de \textit{schema matching}; El Bayadh et al.\ (2015), Chandrasekaran y Mago (2022) y
        Santos Sousa et al.\ (2026) cubren alineación ontológica y similitud semántica. Son técnicas
        aplicables por separado a componentes del contrato OpenAPI y del \ServiceDomain{} BIAN, pero
        ninguna las integra en un score único sobre el par bancario.
  \item \textbf{Separación dominio--interfaz.} Schmidt y Thiry (2020), en una SLR de 27 estudios
        sobre identificación de microservicios con MDE y DDD, muestran cómo las capacidades de
        dominio se traducen en límites de servicio. Ese marco justifica contrastar modelos de negocio
        e interfaces técnicas sin confundir la capa REST con el \ServiceDomain{} BIAN completo.
  \item \textbf{Delimitación del foco.} Las revisiones de testing, seguridad y selección de servicios
        (ejes E5 y E6) emplean OpenAPI con fines operativos ---cobertura de código, vulnerabilidades,
        ranking por QoS--- distintos a la alineación documental; acotan el alcance del presente plan
        (véase \Cref{sec:alcances}).
\end{enumerate}

De estas cuatro líneas se desprende que existen bases parciales para medir correspondencia
estructural (\StructuralScore{}), similitud semántica (\SemanticScore{}) y cobertura conceptual
(\CoverageScore{}), pero \textbf{no hay un procedimiento integrado} que los combine en un
\AlignmentScore{} reproducible sobre contratos OpenAPI bancarios frente a \ServiceDomain{}s BIAN. La
revisión detallada por ejes, las tablas sistematizadas y la formulación explícita de la brecha se
desarrollan a continuación en este capítulo.


\section{Bases conceptuales}
\label{sec:bases-conceptuales}
Las bases conceptuales fijan el vocabulario del dominio bancario--técnico de este plan de tesis.
Se complementan con las definiciones operativas del Capítulo~1 (§1.1 y §1.4) y
con el mapa semántico unificado de la sección~\ref{sec:glosario}.

\subsection{API REST y contrato OpenAPI}

En este plan, el \textbf{contrato OpenAPI} (YAML/JSON) es el artefacto concreto que el banco
publica para describir su API: \textit{paths}, operaciones, \textit{schemas} y metadatos legibles
por herramientas de documentación y análisis. En banca suele documentar pagos, cuentas o préstamos
\textbf{antes} de exponer el servicio a canales digitales o socios \textit{fintech}.

Casas et al.\ (2021) confirman OpenAPI como estándar de facto en APIs web. La capa REST subyacente
---recursos, verbos HTTP y representaciones serializadas en \textit{schemas}--- se toma como
\textbf{contexto}, no como objeto de estudio: lo que se evalúa es el contrato publicado, no la
arquitectura de despliegue ni el tráfico en producción.

\subsection{Modelo de referencia BIAN}

\textbf{BIAN} organiza capacidades bancarias en \ServiceDomain{}s: unidades de negocio descritas
mediante \Behavior{}s (acciones), \BusinessObject{}s (entidades) y relaciones de control. En este
plan, BIAN actúa como \textbf{modelo de referencia} para contrastar si un contrato OpenAPI bancario
refleja la semántica y la estructura esperadas del dominio (Farzi, 2021). No se evalúan
implementaciones en código ni despliegues productivos.

\subsection{Alineación, coherencia y desalineación}

\textbf{Alineación} designa el grado de correspondencia estructural y semántica entre un contrato
OpenAPI bancario y un \ServiceDomain{} BIAN emparejado. \textbf{Coherencia} se emplea como sinónimo
operativo cuando el \AlignmentScore{} integrado alcanza niveles Altos o Medios según la escala del
modelo. \textbf{Desalineación} agrupa las divergencias detectables; la tipología (estructural,
semántica, de modelado, REST/contrato) se define en el Capítulo~1 (Tabla~\ref{tab:tipologias}).

\subsection{Modelo de puntuación}

El \textbf{modelo de alineación} propuesto descompone la evaluación en tres componentes ---
\SemanticScore{} (S), \StructuralScore{} (E) y \CoverageScore{} (C) --- agregados en
\AlignmentScore{} mediante pesos $\alpha$, $\beta$ y $\gamma$. La \Cref{fig:bases-conceptuales}
resume la relación entre capa técnica (OpenAPI), capa de negocio (BIAN) y producto del modelo.

\begin{figure}[H]
  \centering
  \includegraphics[width=0.92\textwidth]{Figura_05_Bases_Conceptuales}
  \caption{Bases conceptuales: contrato OpenAPI, \ServiceDomain{} BIAN y modelo de alineación S/E/C}
  \label{fig:bases-conceptuales}
\end{figure}


\section{Definición de términos (mapa semántico)}
\label{sec:glosario}
El mapa semántico unifica términos empleados en los capítulos del plan. Cuando un concepto aparece
en varios capítulos, conserva el significado aquí definido (PI-01: hilo conductor).

\begin{table}[H]
  \centering
  \caption{Mapa semántico --- términos clave del plan de tesis}
  \label{tab:glosario}
  \small
  \begin{tabular}{@{}p{0.22\textwidth}p{0.68\textwidth}@{}}
    \toprule
    \textbf{Término} & \textbf{Definición en este plan} \\
    \midrule
    API REST &
    Interfaz que expone recursos y operaciones HTTP con representaciones intercambiables (p.\ ej. JSON). \\
    Contrato OpenAPI &
    Documento YAML/JSON (OpenAPI~3.x) que describe \textit{paths}, \textit{operations} y \textit{schemas} de una API. \\
    BIAN &
    Estándar de arquitectura bancaria que organiza capacidades en \ServiceDomain{}s de referencia. \\
    \ServiceDomain{} &
    Unidad de negocio BIAN descrita por \Behavior{}s, \BusinessObject{}s y relaciones de control. \\
    \Behavior{} &
    Acción o capacidad de negocio definida en un \ServiceDomain{} BIAN. \\
    \BusinessObject{} &
    Entidad de dominio BIAN (p.\ ej. orden de pago, cuenta) referenciada en el modelo. \\
    Objeto de estudio &
    Contrato OpenAPI bancario completo evaluado frente a un \ServiceDomain{} BIAN. \\
    Unidad de análisis &
    Par comparable endpoint/operation/schema $\leftrightarrow$ \Behavior{} / \BusinessObject{}. \\
    Alineación &
    Grado de correspondencia estructural y semántica entre contrato OpenAPI y \ServiceDomain{} BIAN. \\
    Desalineación &
    Divergencia detectable; tipologías: estructural, semántica, modelado, REST/contrato (Tabla~\ref{tab:tipologias}). \\
    Representación normalizada &
    Vista intermedia común (tabular, grafo u ontología) de artefactos OpenAPI y BIAN tras OE1. \\
    \textit{Schema matching} &
    Técnicas de correspondencia entre esquemas o estructuras; base del \StructuralScore{} (E). \\
    Similitud semántica &
    Medida de equivalencia de significado entre términos/conceptos; base del \SemanticScore{} (S). \\
    \SemanticScore{} (S) &
    Score $\in [0,1]$ de alineación semántica OpenAPI $\leftrightarrow$ BIAN. \\
    \StructuralScore{} (E) &
    Score $\in [0,1]$ de correspondencia estructural REST $\leftrightarrow$ organización del SD. \\
    \CoverageScore{} (C) &
    Score $\in [0,1]$ de conceptos BIAN cubiertos por el contrato OpenAPI. \\
    \AlignmentScore{} &
    Agregación $\alpha S + \beta E + \gamma C$; clasificación Alta/Media/Baja/Nula. \\
    \textit{Matching} &
    Procedimiento que establece correspondencias entre elementos de dos artefactos normalizados. \\
    \bottomrule
  \end{tabular}
\end{table}


Los términos \SemanticScore{}, \StructuralScore{}, \CoverageScore{} y \AlignmentScore{} se usan
con la convención \textbf{S}, \textbf{E}, \textbf{C} del modelo (véase §1.4 y Anexo~C). El
\StructuralScore{} mide correspondencia REST $\leftrightarrow$ organización del \ServiceDomain{};
no debe confundirse con similitud ontológica genérica ni con métricas de \textit{runtime}.


\section{Bases teóricas}
\label{sec:bases-teoricas}
Las bases teóricas sustentan las técnicas que operacionalizan los objetivos específicos del
Capítulo~2 y el diseño metodológico del Capítulo~4. Se apoyan en la literatura revisada en
§\ref{sec:estado-arte}; los papers de \textit{aporte} citados aquí fundan \textbf{conceptos} de la
solución propuesta, no sustituyen al estado del arte.

\subsection{Correspondencia estructural (\textit{schema matching})}

Shvaiko y Euzenat (2005) proponen una taxonomía multidimensional del \textbf{\textit{schema
matching}}: técnicas aproximadas o exactas a nivel de esquema; evidencia sintáctica, semántica o
externa a nivel de elemento y estructura; grado de automatización. El marco es anterior a OpenAPI y
JSON Schema, pero fundamenta el cálculo del \StructuralScore{} al contrastar \textit{paths},
\textit{operations} y \textit{schemas} con la organización de \Behavior{}s y \BusinessObject{}s del
\ServiceDomain{} de referencia (ficha E2-01, Anexo~B).

\subsection{Alineación ontológica y similitud semántica}

El Bayadh et al.\ (2015) clasifican enfoques de \textbf{alineación ontológica} en cuatro familias:
medidas terminológicas, estructurales, extensionales y semánticas. Ninguna medida aislada basta;
combinaciones fijas pueden ser subóptimas, lo que justifica un diseño explícito de pesos en el
modelo S/E/C del Grupo~9.2.

Chandrasekaran y Mago (2022) trazan la evolución de la \textbf{similitud semántica} desde recursos
léxicos (WordNet) hacia \textit{embeddings} y modelos transformadores. Santos Sousa et al.\ (2026)
catalogan el uso de \textit{embeddings} en ontology matching, señalando limitaciones de supervisión,
escalabilidad y correspondencias complejas.

Estas líneas aportan candidatos para el \SemanticScore{} y, en parte, para estimar cobertura
conceptual (\CoverageScore{}), siempre sobre vocabulario extraído del contrato OpenAPI y del modelo
BIAN normalizado.

\subsection{Representación enriquecida de contratos}

Mainas et al.\ (2023) ---citados en el eje E1--- proponen anotaciones semánticas
(\texttt{x-refersTo}, \texttt{x-kindOf}, etc.), ontología REST/Hydra y traducción OpenAPI
$\rightarrow$ RDF como vía de \textbf{representación enriquecida}. Evaluaron más de 300 contratos de
SwaggerHub con consultas SPARQL; los tiempos de traducción superan 1~s para YAML de más de 200
líneas. En este plan se consideran como técnicas de normalización candidatas (Cap.~4), no como
definición previa del \SemanticScore{}.

\subsection{Relación con el modelo propuesto}

Las bases teóricas no constituyen por sí solas el producto de la investigación: integran técnicas
parciales (E2, E3) en un \textbf{procedimiento de diseño} que produce S, E, C y \AlignmentScore{}
sobre artefactos documentales bancarios. La brecha que justifica esa integración se expone en
§\ref{sec:brecha}.


\section{Estado del arte por ejes temáticos}
\label{sec:estado-arte}
El estado del arte se organiza en \textbf{seis ejes} (E1--E6) del Anexo~B. Cada eje responde a una
pregunta sobre OpenAPI, \textit{matching} estructural o semántico, arquitectura de servicios o
delimitación del foco. La \Cref{tab:sintesis-ejes} resume hallazgos convergentes; el Anexo~B detalla
cada paper.

\subsection{Eje E1 --- OpenAPI como estándar}

\textbf{Papers:} Casas et al.\ (2021), Mainas et al.\ (2023), Santos y Casas (2025).

Casas et al.\ (2021) realizan un mapeo sistemático de 47 artículos primarios (2011--2020) sobre
usos de OpenAPI/Swagger. Confirman que es el \textbf{estándar de facto} para describir APIs REST de
forma \textit{machine-readable}: el 68\,\% de los trabajos mapeados propone herramientas o
automatización; el 43\,\% mejora documentación. Sin embargo, ningún estudio del eje vincula esos
contratos con modelos de referencia de negocio bancario tipo BIAN.

Mainas et al.\ (2023) avanzan hacia \textbf{OpenAPI enriquecido semánticamente}. Proponen
anotaciones (\texttt{x-refersTo}, \texttt{x-kindOf}, \texttt{x-mapsTo}, etc.), ontología REST/Hydra,
traducción OpenAPI $\rightarrow$ RDF y consultas SPARQL sobre un corpus de más de 300 contratos de
SwaggerHub. Aportan \textbf{técnicas de representación y normalización} candidatas para calcular
\SemanticScore{}; no definen S ni \AlignmentScore{}, y operan en dominio genérico sin
\ServiceDomain{} BIAN.

Santos y Casas (2025), mediante Design Science Research, construyen un catálogo de 59 métricas de
\textbf{API Gateway} (latencia, tráfico, errores) evaluadas en 68 productos comerciales. Delimitan un
borde útil para este estudio: la calidad operativa en producción \textbf{no equivale} a la coherencia
documental OpenAPI $\leftrightarrow$ modelo de negocio.

\textbf{Conclusión E1:} OpenAPI está maduro como artefacto técnico; la literatura no resuelve su
\textbf{alineación} con capacidades de negocio estandarizadas en banca.

\subsection{Eje E2 --- Matching estructural}

\textbf{Paper:} Shvaiko y Euzenat (2005).

Shvaiko y Euzenat (2005) sistematizan el \textbf{\textit{schema matching}} en una taxonomía que
cruza aproximación exacta/approximada, evidencia sintáctica/semántica/externa y nivel
elemento/estructura. El marco es anterior a OpenAPI y JSON Schema ---no aborda contratos REST ni
dominios bancarios---, pero fundamenta el \StructuralScore{} al contrastar \textit{paths},
\textit{operations} y \textit{schemas} con la organización de un \ServiceDomain{} BIAN
(\Behavior{}, \BusinessObject{}) como problema de correspondencia entre esquemas.

\textbf{Conclusión E2:} Existe base teórica sólida para correspondencias estructurales; falta
aplicarla a \textbf{pares OpenAPI bancario $\leftrightarrow$ BIAN} con métrica reproducible.

\subsection{Eje E3 --- Matching semántico y ontologías}

\textbf{Papers:} El Bayadh et al.\ (2015), Chandrasekaran y Mago (2022), Santos Sousa et al.\ (2026).

Tres líneas convergen. El Bayadh et al.\ (2015) clasifican \textbf{alineación ontológica} en medidas
terminológicas, estructurales, extensionales y semánticas; advierten que ninguna medida aislada basta.
Chandrasekaran y Mago (2022) trazan la evolución de \textbf{similitud semántica}
desde recursos léxicos hasta \textit{embeddings} y \textit{transformers}. Santos Sousa
et al.\ (2026) revisan \textbf{\textit{embeddings} para ontology matching}, con limitaciones de
supervisión, escalabilidad y correspondencias complejas.

Ninguno aborda OpenAPI ni BIAN directamente; todos aportan técnicas candidatas para
\SemanticScore{} y, en parte, cobertura conceptual (\CoverageScore{}).

\textbf{Conclusión E3:} La similitud semántica está bien documentada; no existe propuesta integrada
que la aplique a \textbf{vocabulario OpenAPI bancario vs.\ conceptos BIAN} dentro de un modelo de
alineación único.

\subsection{Eje E4 --- Microservicios e interfaces}

\textbf{Paper:} Schmidt y Thiry (2020).

Schmidt y Thiry (2020) revisan 27 estudios primarios (de 715 registros, 2013--2019) sobre
estrategias de \textbf{identificación de microservicios} con MDE y DDD. La síntesis muestra cómo la
descomposición desde modelos de dominio hacia límites de servicio contextualiza por qué los
\ServiceDomain{} BIAN son referencia de negocio y las APIs su capa de exposición. Pocos trabajos
integran MDE y DDD conjuntamente; ninguno mide coherencia entre contrato OpenAPI y modelo BIAN.

\textbf{Conclusión E4:} La literatura arquitectónica relaciona dominio e interfaces; \textbf{no
cuantifica} desalineación OpenAPI--BIAN.

\subsection{Eje E5 --- Testing y calidad API (borde del estudio)}

\textbf{Papers:} Ehsan et al.\ (2022), Golmohammadi et al.\ (2024), Otieno y Araka (2026); Tanveer
et al.\ (2025), Vitorino et al.\ (2025).

Cinco revisiones sistemáticas cubren \textbf{testing REST}, ML para pruebas, \textbf{vulnerabilidades}
y generación adversarial de datos. En todos los casos OpenAPI entra como \textbf{entrada} a pruebas
automatizadas, no como artefacto a contrastar con un \ServiceDomain{} BIAN.

Ehsan et al.\ (2022) incluyen 16 artículos sobre metodologías de testing RESTful. La
\textit{code coverage} es la métrica principal de efectividad.

Golmohammadi et al.\ (2024) analizan 92 papers con taxonomía de técnicas (\textit{fuzzing},
\textit{search-based}, model-based). La mayoría consume contratos OpenAPI para generar casos de
prueba sobre endpoints ya desplegados.

Otieno y Araka (2026) revisan ML aplicado a testing RESTful; reportan datasets pequeños y costo de
etiquetado como limitaciones recurrentes.

Tanveer et al.\ (2025) sistematizan detección de vulnerabilidades en APIs REST; el foco es seguridad
\textit{runtime}, ajeno al \AlignmentScore{} documental.

Vitorino et al.\ (2025) abordan aprendizaje adversarial para pruebas automatizadas. Confirman que la
literatura de calidad operativa \textbf{no sustituye} un procedimiento de alineación OpenAPI--BIAN.

\textbf{Conclusión E5:} Confirma la delimitación del plan (\Cref{sec:alcances}): testing, seguridad
y QoS operativa \textbf{quedan fuera}; el producto central es el \AlignmentScore{} documental.

\subsection{Eje E6 --- Selección y recomendación de servicios}

\textbf{Paper:} Mahanra Rao et al.\ (2025).

Mahanra Rao et al.\ (2025) revisan 25 estudios (2020--2024, PRISMA) sobre enfoques
\textbf{BERT/transformers} para selección y recomendación de servicios web con foco en predicción de
QoS. La representación semántica de descripciones de servicio es adyacente a S, pero el objetivo es
\textbf{ranking/consumo}, no validar un contrato frente a un \ServiceDomain{} de referencia. Los
autores señalan alto costo computacional, datasets limitados e interpretabilidad reducida.

\textbf{Conclusión E6:} Contexto semántico de servicios; \textbf{no sustituye} un modelo de
alineación OpenAPI $\leftrightarrow$ BIAN.

\begin{table}[H]
  \centering
  \caption{Síntesis del estado del arte por ejes (E1--E6)}
  \label{tab:sintesis-ejes}
  \small
  \begin{tabular}{@{}p{0.06\textwidth}p{0.20\textwidth}p{0.22\textwidth}p{0.16\textwidth}p{0.24\textwidth}@{}}
    \toprule
    \textbf{Eje} & \textbf{Papers} & \textbf{Hallazgo principal} & \textbf{Aporte S/E/C} & \textbf{Vacío OpenAPI--BIAN} \\
    \midrule
    \textbf{E1} &
    Casas 2021; Mainas 2023; Santos y Casas 2025 &
    OpenAPI consolidado; enriquecimiento semántico posible &
    S (Mainas); contexto estándar &
    Sin vínculo con \ServiceDomain{} BIAN \\
    \textbf{E2} &
    Shvaiko 2005 &
    Taxonomía de \textit{schema matching} &
    E &
    Sin aplicación REST/banca/BIAN \\
    \textbf{E3} &
    El Bayadh 2015; Chandrasekaran 2022; Santos Sousa 2026 &
    Ontología matching + embeddings + STS &
    S; C parcial &
    Sin contratos OpenAPI ni dominio bancario \\
    \textbf{E4} &
    Schmidt 2020 &
    Dominio (DDD/MDE) $\rightarrow$ microservicios &
    Contexto dominio$\rightarrow$interfaz (no \StructuralScore{} operativo) &
    Sin métrica de alineación \\
    \textbf{E5} &
    Ehsan 2022; Golmohammadi 2024; Otieno 2026; Tanveer 2025; Vitorino 2025 &
    Testing, ML-test, seguridad &
    Delimitación (fuera de foco) &
    Orthogonal al \AlignmentScore{} \\
    \textbf{E6} &
    Mahanra Rao 2025 &
    BERT + QoS en selección de servicios &
    S parcial &
    Recomendación $\neq$ alineación BIAN \\
    \bottomrule
  \end{tabular}
\end{table}


La Tabla~7 (Anexo~B) detalla objetivo, técnicas, métricas y limitaciones \textbf{paper por paper}.


\section{Brecha y oportunidad de investigación}
\label{sec:brecha}
\subsection{Qué sí cubre la literatura revisada}

Esta subsección sintetiza, de forma transversal, lo que los catorce papers aportan al diseño del
modelo. No sustituye la lectura por ejes en §\ref{sec:estado-arte}.

\begin{enumerate}
  \item \textbf{OpenAPI} es el formato dominante para especificar APIs REST (E1). Casas et al.\
        (2021) mapean 47 artículos que lo confirman; Mainas et al.\ (2023) muestran enriquecimiento
        semántico viable sobre corpus reales.
  \item Existen marcos para \textbf{correspondencia entre esquemas} (E2) y \textbf{similitud
        semántica} (E3). Shvaiko y Euzenat (2005) y El Bayadh et al.\ (2015) ofrecen taxonomías
        reutilizables; Santos Sousa et al.\ (2026) documentan limitaciones prácticas de
        \textit{embeddings}.
  \item La \textbf{arquitectura orientada al dominio} separa modelos de negocio e interfaces
        técnicas (E4). Schmidt y Thiry (2020) relacionan capacidades de dominio con microservicios,
        en la misma línea que BIAN como referencia.
  \item Por el contrario, la literatura de \textbf{testing y seguridad} (E5) es abundante: Ehsan
        et al.\ (2022) y Golmohammadi et al.\ (2024) usan OpenAPI como entrada a pruebas
        automatizadas. Esa línea ayuda a \textbf{delimitar} el plan ---calidad operativa en
        producción no equivale a coherencia documental con BIAN.
  \item En selección de servicios (E6), Mahanra Rao et al.\ (2025) optimizan ranking por QoS con
        \textit{embeddings}; el objetivo es consumo, no verificar si un contrato bancario refleja un
        \ServiceDomain{} de referencia.
\end{enumerate}

Ninguna línea, sola o en conjunto, responde a la pregunta del plan: \textit{¿en qué medida un
contrato OpenAPI bancario refleja el modelo de un \ServiceDomain{} BIAN y cómo cuantificarlo de
forma reproducible?}

\subsection{Vacíos identificados}

\begin{table}[H]
  \centering
  \caption{Vacíos identificados en la literatura (V1--V5)}
  \label{tab:vacios}
  \small
  \begin{tabular}{@{}p{0.14\textwidth}p{0.38\textwidth}p{0.38\textwidth}@{}}
    \toprule
    \textbf{Vacío} & \textbf{Evidencia en el review} & \textbf{Implicación para la tesis} \\
    \midrule
    \textbf{V1 --- Intersección OpenAPI $\cap$ BIAN} &
    Ningún paper estudia ambos artefactos en el sector bancario &
    El objeto de estudio (contrato OpenAPI bancario vs.\ SD BIAN) no está cubierto \\
    \textbf{V2 --- Métrica integrada} &
    E2 y E3 aportan técnicas parciales; E1 no evalúa alineación &
    Justifica S, E, C y \AlignmentScore{} como producto original \\
    \textbf{V3 --- Cobertura de conceptos} &
    Matching ontológico discute cobertura; no sobre \BusinessObject{} BIAN en OpenAPI &
    Justifica \CoverageScore{} (C) \\
    \textbf{V4 --- Procedimiento reproducible} &
    Mainas propone \textit{pipeline} RDF; no scores comparables ni clasificación &
    Justifica diseño metodológico Cap.~4 (Guía~1) \\
    \textbf{V5 --- Gobernanza documental} &
    E5/E6 abordan \textit{runtime} o selección, no deriva OpenAPI--modelo de referencia &
    Motiva alineación previa a despliegue/certificación \\
    \bottomrule
  \end{tabular}
\end{table}


La \Cref{fig:brecha} sintetiza la transición desde los ejes del estado del arte hacia los vacíos
V1--V5 y la oportunidad de investigación.

\begin{figure}[H]
  \centering
  \includegraphics[width=0.95\textwidth]{Figura_06_Brecha}
  \caption{De la literatura revisada (E1--E6) a la brecha de investigación (V1--V5)}
  \label{fig:brecha}
\end{figure}

\subsection{Oportunidad de investigación}

La brecha abre la oportunidad de \textbf{diseñar un modelo de alineación} documentado en cinco
fases (Fase~1: diseño; Fases~2--5: transformaciones OE1--OE4). Dado un contrato OpenAPI bancario y
un \ServiceDomain{} BIAN emparejado, el modelo debe:

\begin{itemize}
  \item normalizar ambos artefactos a representaciones comparables (OE1);
  \item calcular \SemanticScore{}, \StructuralScore{} y \CoverageScore{} (OE2--OE3);
  \item integrar \AlignmentScore{} $= \alpha \cdot S + \beta \cdot E + \gamma \cdot C$ con
        clasificación interpretativa (Alta / Media / Baja / Nula).
\end{itemize}

Este plan documenta el \textbf{diseño} del procedimiento. La construcción del artefacto y la
experimentación empírica quedan fuera de alcance (\Cref{sec:alcances}).

La investigación contribuye a los vacíos V1--V4 con un enfoque reproducible sobre artefactos
documentales. No sustituye testing operativo (E5) ni certificación regulatoria; tampoco confunde
calidad \textit{runtime} (E1/E6) con coherencia documental (V5).


En conjunto, el marco teórico delimita la brecha V1--V5 y prepara las técnicas del
Capítulo~\ref{cap:metodologia}, manteniendo separación entre estado del arte y problema tecnológico
de la matriz (Anexo~C).
